% GENERATED FILE. Edit canonical lessons, then run npm run generate:books.
% canonical-source-hash: fnv1a64:bbce5f1707b5e6b7
% canonical-lessons: ML-C366-sucana, ML-C366-atayalam, ML-C366-kalppatu, ML-C366-viralatayalam, ML-C366-duruhamaya

\chapter{Clues and Footprints}
\label{ch:ml-clues-and-footprints}
\hlchaptermodality{366}

\begin{chapteropening}
\textbf{By the end of this chapter:} \emph{I can say a clue, a mark, a footprint, a fingerprint and mysterious.}

Complete the last lesson of chapter 366: I can say a clue, a mark, a footprint, a fingerprint and mysterious.
\end{chapteropening}

\section[\texorpdfstring{\ml{സൂചന} (sūcana)}{സൂചന (sūcana)}]{\ml{സൂചന} (sūcana) — a clue; a hint}

\label{lesson:ML-C366-sucana}



\begin{quote}
Before the new one: say the Malayalam for to arrest, then the Malayalam for to seize.
\end{quote}

\subsection*{You'll want to know: \ml{സൂചന}}
\textbf{\ml{സൂചന}} — \emph{sūcana} — \textquotedblleft{}a clue; a hint\textquotedblright{}.

\begin{quote}
\ml{പോലീസിന്} \ml{കള്ളനെക്കുറിച്ച്} \ml{ഒരു} \ml{സൂചന} \ml{കിട്ടി}.
\end{quote}

\begin{quote}
\emph{pōlīsinŭ kaḷḷanekkuṟiccŭ oru sūcana kiṭṭi.}
\end{quote}

\begin{quote}
The police got a clue about the thief.
\end{quote}

Also a hint or a sign: \ml{മഴയുടെ} \ml{സൂചന}, a sign of rain.

\begin{cousinweb}
\textbf{\ml{സൂചന}} (\emph{sūcana}) is borrowed from Sanskrit \emph{sūcanā}, \textquotedblleft{}pointing out, indication, communication\textquotedblright{} (Monier-Williams).
\end{cousinweb}

\begin{grammarlens}[title={What you've built}]
One more word for telling what happened.
\end{grammarlens}

\subsection*{Your turn}
\begin{itemize}
\raggedright
  \item \emph{Say it:} \emph{sūcana}
  \item \emph{Say it:} \emph{sūcana}, once more
  \item \emph{Say it:} \emph{pōlīsinŭ kaḷḷanekkuṟiccŭ oru sūcana kiṭṭi.} — The police got a clue about the thief.
  \item \emph{Recall:} say the Malayalam for to arrest, then the Malayalam for to seize, then say \emph{sūcana} again
\end{itemize}

\begin{tcolorbox}[breakable,colback=teal!4,colframe=teal!35!black,title={Before you move on}]
What does \ml{സൂചന} mean? (\textquotedblleft{}A clue; a hint\textquotedblright{}.) Say it once more.
\end{tcolorbox}
\section[\texorpdfstring{\ml{അടയാളം} (aṭayāḷaṁ)}{അടയാളം (aṭayāḷaṁ)}]{\ml{അടയാളം} (aṭayāḷaṁ) — a mark; a sign}

\label{lesson:ML-C366-atayalam}



\begin{quote}
Before the new one: say the Malayalam for to seize, then the Malayalam for a clue.
\end{quote}

\subsection*{You'll want to know: \ml{അടയാളം}}
\textbf{\ml{അടയാളം}} — \emph{aṭayāḷaṁ} — \textquotedblleft{}a mark; a sign\textquotedblright{}.

\begin{quote}
\ml{കള്ളൻ} \ml{ഭിത്തിയിൽ} \ml{ഒരു} \ml{അടയാളം} \ml{ഇട്ടു}.
\end{quote}

\begin{quote}
\emph{kaḷḷan bhittiyil oru aṭayāḷaṁ iṭṭu.}
\end{quote}

\begin{quote}
The thief left a mark on the wall.
\end{quote}

Also a symbol or a landmark.

\begin{cousinweb}
\textbf{\ml{അടയാളം}} (\emph{aṭayāḷaṁ}) is native Dravidian: DEDR 89 lists Malayalam \emph{aṭayāḷam}, \textquotedblleft{}mark, sign\textquotedblright{}, beside Tamil \emph{aṭaiyāḷam}, \textquotedblleft{}mark, symbol, seal\textquotedblright{}.
\end{cousinweb}

\begin{grammarlens}[title={What you've built}]
One more word for telling what happened.
\end{grammarlens}

\subsection*{Your turn}
\begin{itemize}
\raggedright
  \item \emph{Say it:} \emph{aṭayāḷaṁ}
  \item \emph{Say it:} \emph{aṭayāḷaṁ}, once more
  \item \emph{Say it:} \emph{kaḷḷan bhittiyil oru aṭayāḷaṁ iṭṭu.} — The thief left a mark on the wall.
  \item \emph{Recall:} say the Malayalam for to seize, then the Malayalam for a clue, then say \emph{aṭayāḷaṁ} again
\end{itemize}

\begin{tcolorbox}[breakable,colback=teal!4,colframe=teal!35!black,title={Before you move on}]
What does \ml{അടയാളം} mean? (\textquotedblleft{}A mark; a sign\textquotedblright{}.) Say it once more.
\end{tcolorbox}
\section[\texorpdfstring{\ml{കാൽപ്പാട്} (kālppāṭŭ)}{കാൽപ്പാട് (kālppāṭŭ)}]{\ml{കാൽപ്പാട്} (kālppāṭŭ) — a footprint}

\label{lesson:ML-C366-kalppatu}



\begin{quote}
Before the new one: say the Malayalam for a clue, then the Malayalam for a mark.
\end{quote}

\subsection*{You'll want to know: \ml{കാൽപ്പാട്}}
\textbf{\ml{കാൽപ്പാട്}} — \emph{kālppāṭŭ} — \textquotedblleft{}a footprint\textquotedblright{}.

\begin{quote}
\ml{മണലിൽ} \ml{കള്ളന്റെ} \ml{കാൽപ്പാടുകൾ} \ml{ഉണ്ടായിരുന്നു}.
\end{quote}

\begin{quote}
\emph{maṇalil kaḷḷanṟe kālppāṭukaḷ uṇṭāyirunnu.}
\end{quote}

\begin{quote}
The thief's footprints were in the sand.
\end{quote}

Plural: \ml{കാൽപ്പാടുകൾ}, \emph{kālppāṭukaḷ}.

\ml{കാൽ}, foot, + \ml{പാട്}, a mark, a trace.

\begin{cousinweb}
\textbf{\ml{കാൽപ്പാട്}} (\emph{kālppāṭŭ}) is a native Dravidian compound of \ml{കാൽ} (\emph{kāl}), \textquotedblleft{}leg, foot\textquotedblright{} (DEDR 1479), and \ml{പാട്}, \textquotedblleft{}a mark, a trace\textquotedblright{}: the mark a foot leaves.
\end{cousinweb}

\begin{grammarlens}[title={What you've built}]
One more word for telling what happened.
\end{grammarlens}

\subsection*{Your turn}
\begin{itemize}
\raggedright
  \item \emph{Say it:} \emph{kālppāṭŭ}
  \item \emph{Say it:} \emph{kālppāṭŭ}, once more
  \item \emph{Say it:} \emph{maṇalil kaḷḷanṟe kālppāṭukaḷ uṇṭāyirunnu.} — The thief's footprints were in the sand.
  \item \emph{Recall:} say the Malayalam for a clue, then the Malayalam for a mark, then say \emph{kālppāṭŭ} again
\end{itemize}

\begin{tcolorbox}[breakable,colback=teal!4,colframe=teal!35!black,title={Before you move on}]
What does \ml{കാൽപ്പാട്} mean? (\textquotedblleft{}A footprint\textquotedblright{}.) Say it once more.
\end{tcolorbox}
\section[\texorpdfstring{\ml{വിരലടയാളം} (viralaṭayāḷaṁ)}{വിരലടയാളം (viralaṭayāḷaṁ)}]{\ml{വിരലടയാളം} (viralaṭayāḷaṁ) — a fingerprint}

\label{lesson:ML-C366-viralatayalam}



\begin{quote}
Before the new one: say the Malayalam for a mark, then the Malayalam for a footprint.
\end{quote}

\subsection*{You'll want to know: \ml{വിരലടയാളം}}
\textbf{\ml{വിരലടയാളം}} — \emph{viralaṭayāḷaṁ} — \textquotedblleft{}a fingerprint\textquotedblright{}.

\begin{quote}
\ml{പോലീസ്} \ml{വാതിലിൽ} \ml{നിന്ന്} \ml{വിരലടയാളം} \ml{എടുത്തു}.
\end{quote}

\begin{quote}
\emph{pōlīsŭ vātilil ninnŭ viralaṭayāḷaṁ eṭuttu.}
\end{quote}

\begin{quote}
The police took fingerprints from the door.
\end{quote}

\ml{വിരൽ}, finger, + \ml{അടയാളം}, a mark.

\begin{cousinweb}
\textbf{\ml{വിരലടയാളം}} (\emph{viralaṭayāḷaṁ}) is a native Dravidian compound of \ml{വിരൽ} (\emph{viral}), \textquotedblleft{}finger\textquotedblright{} (DEDR 5409), and \ml{അടയാളം} (\emph{aṭayāḷaṁ}), \textquotedblleft{}a mark\textquotedblright{} (DEDR 89).
\end{cousinweb}

\begin{grammarlens}[title={What you've built}]
One more word for telling what happened.
\end{grammarlens}

\subsection*{Your turn}
\begin{itemize}
\raggedright
  \item \emph{Say it:} \emph{viralaṭayāḷaṁ}
  \item \emph{Say it:} \emph{viralaṭayāḷaṁ}, once more
  \item \emph{Say it:} \emph{pōlīsŭ vātilil ninnŭ viralaṭayāḷaṁ eṭuttu.} — The police took fingerprints from the door.
  \item \emph{Recall:} say the Malayalam for a mark, then the Malayalam for a footprint, then say \emph{viralaṭayāḷaṁ} again
\end{itemize}

\begin{tcolorbox}[breakable,colback=teal!4,colframe=teal!35!black,title={Before you move on}]
What does \ml{വിരലടയാളം} mean? (\textquotedblleft{}A fingerprint\textquotedblright{}.) Say it once more.
\end{tcolorbox}
\section[\texorpdfstring{\ml{ദുരൂഹമായ} (durūhamāya)}{ദുരൂഹമായ (durūhamāya)}]{\ml{ദുരൂഹമായ} (durūhamāya) — mysterious}

\label{lesson:ML-C366-duruhamaya}



\begin{quote}
Before the new one: say the Malayalam for a footprint, then the Malayalam for a fingerprint.
\end{quote}

\subsection*{You'll want to know: \ml{ദുരൂഹമായ}}
\textbf{\ml{ദുരൂഹമായ}} — \emph{durūhamāya} — \textquotedblleft{}mysterious\textquotedblright{}.

\begin{quote}
\ml{ആ} \ml{കത്ത്} \ml{വളരെ} \ml{ദുരൂഹമായിരുന്നു}.
\end{quote}

\begin{quote}
\emph{ā kattŭ vaḷare durūhamāyirunnu.}
\end{quote}

\begin{quote}
That letter was very mysterious.
\end{quote}

The noun is \ml{ദുരൂഹത}, a mystery.

\begin{cousinweb}
\textbf{\ml{ദുരൂഹമായ}} (\emph{durūhamāya}) is built from Sanskrit \emph{durūha}, \textquotedblleft{}hard to infer or understand\textquotedblright{} (Monier-Williams), with the Malayalam adjective ending -\ml{ആയ}: a mysterious thing is one nobody can work out.
\end{cousinweb}

\begin{grammarlens}[title={What you've built}]
One more word for telling what happened.
\end{grammarlens}

\subsection*{Your turn}
\begin{itemize}
\raggedright
  \item \emph{Say it:} \emph{durūhamāya}
  \item \emph{Say it:} \emph{durūhamāya}, once more
  \item \emph{Say it:} \emph{ā kattŭ vaḷare durūhamāyirunnu.} — That letter was very mysterious.
  \item \emph{Recall:} say the Malayalam for a footprint, then the Malayalam for a fingerprint, then say \emph{durūhamāya} again
\end{itemize}

\begin{tcolorbox}[breakable,colback=teal!4,colframe=teal!35!black,title={Before you move on}]
What does \ml{ദുരൂഹമായ} mean? (\textquotedblleft{}Mysterious\textquotedblright{}.) Say it once more.
\end{tcolorbox}
